% ===================== 第 12 章 =====================
\section{设计边界与演进路线}

一套以审计为核心价值的系统，必须能准确说出自己这一版覆盖到哪里、没覆盖到哪里。本章把分散在正文的范围界定集中陈述。以下每条都是**设计决策的产物**——首版围绕「组织结构能否提升盲点发现质量」这一核心主张收敛范围，凡是与该主张无关的能力一律后置（CLAUDE.md 第 14 条 YAGNI）。边界不等于欠账，它划定了当前版本验证结论的适用域。

\subsection{四项版本边界}

\begin{enumerate}[leftmargin=1.6em, itemsep=4pt]
\item \textbf{记忆层以冻结接口接入，首版实现为内存适配器。}三层接口（\texttt{init\_private\_memory} / \texttt{memorize\_episode} / \texttt{recall\_private} / \texttt{save\_snapshot} / \texttt{load\_snapshot}）与席位隔离规则已经冻结，替换为上游 MemoBrain 实现不需要改动议会层与证据层的任何代码。首版运行在内存适配器与简化折叠实现之上，因此「长程记忆防止漂移」这一收益当前由结构保证承载，其跨任务实测数据属于下一阶段。
\item \textbf{证据分级采取保守上限：B 级为独立支撑上限。}首版取证管线以题录与摘要为稳定供给面，全文获取与独立双抽取作为进阶通道实现，覆盖面仍在扩展。系统因此主动禁止 B 级证据单独支撑高置信因果结论。这是刻意的保守选择：宁可让结论的证据强度上限明确可见，也不让一条 A 级才配得上的因果表述借 B 级材料过关。
\item \textbf{完整协议的算力需求与模型档位绑定。}第 10 章真实模型实验显示，flash 级模型在完整议会下产生大量空槽（12.3 个未填充证据槽位，简化变体为 1.0--1.7），治理机制「以召回换精度」的收益要在推理级模型或更轻的协议变体上才能被公平观测。这不是协议本身的缺陷，而是协议与推理预算之间的匹配问题；系统按设计报告自己正运行在权衡的哪一侧，而不是把两侧平均成一个分数。
\item \textbf{评测语料与人工金标准为独立的并行工程。}ForesightBlindspot 的封闭语料建设与双人标注一致性数据与系统实现并行推进。在该语料交付之前，评测以「脚本化受控案例 + 真实模型重复运行 + 语义裁判」为口径；所有评测数字必须连同其口径（脚本化 / 真实模型、重复次数、何种裁判）一起阅读。
\end{enumerate}

\subsection{明确不做的事（YAGNI）}

首版范围刻意排除：通用全学科、自动实验设备、任意代码复现、自动元分析、模型训练平台、微服务拆分、多租户复杂权限、插件市场、移动端。这些方向对证明「组织结构提升盲点发现质量」这一首版主张没有贡献，纳入只会稀释验证强度。

\subsection{演进路线（按对核心主张的贡献排序）}

\begin{table}[ht]
\centering\small
\begin{tabularx}{\linewidth}{@{}L{1.6cm} L{3.6cm} X@{}}
\toprule
\rowcolor{accent}
\tblhead 优先级 & \tblhead 方向 & \tblhead 要回答的问题 \\
\midrule
P0 & 推理级模型 / 轻协议对照实验 & 在模型能稳定产出结构化动作的条件下，完整议会的精度—召回曲线是否稳定优于消融变体 \\
P0 & 时间切片语料建设 & 封闭语料 + 双人标注金标准，产出独立 B\_test 与一致性系数 \\
P1 & 上游 MemoBrain 真实接入 & 跨任务长程记忆对盲点召回与长程漂移的实测影响 \\
P1 & 全文获取与独立双抽取 & 把可支撑结论的证据等级从 B 推进到 A \\
P2 & 盲证据评审轮次处理器 & 把当前的结构性约束升级为主动剥离声誉字段的完整实现 \\
P2 & 证据图的研究者编辑工作流 & 在不绕过证据门的前提下，支持研究者对主张/边界做受控标注与导出 \\
\bottomrule
\end{tabularx}
\end{table}

\subsection{结语}

Poliscope 的赌注很具体：当多智能体研究系统的讨论从「谁的总结更流畅」转向「哪条结论被什么级别的证据支撑、谁反对、需要什么才能改判」，科研审查才真正发生。七人议会、事件账本、双图隔离与时间切片评测，都是为了让这三个问题在系统的每一层都有确定的答案——包括在系统暂时回答不了的时候，明确地把「不知道」显示出来。
